\section{Discussion}\label{sec:discussion}

The clearest result in this paper is not a broad improvement in technology transfer after universities revise their IP policies. The pattern is much narrower. Licensing follows different paths after upward and downward revisions, while the evidence for earlier stages of the commercialization process is considerably weaker. The filing margin and new patent applications move in the same direction as licensing but are estimated less precisely, and invention disclosures do not rise. The licensing result also becomes less compelling once we adjust for examining several outcomes. We therefore see the evidence as pointing to a particular part of the commercialization process rather than to a general effect of policy revision.

This matters for how our findings relate to the earlier literature. Much of the work on university IP policy has focused on ownership and financial incentives. That is a natural place to begin: ownership determines who controls an invention, and royalty sharing rules change the financial return to the inventor \citep{HvideJones2018,KenneyPatton2009,LachSchankerman2008,ArqueCastellsEtAl2016}. But the evidence on these incentives has not always been stable. \citet{OuelletteTutt2020}, for example, show that the earlier relationship between royalty shares and licensing income was sensitive to errors in the underlying policy data. Our setting looks at a different margin. The revisions we study usually take place within an existing system of university ownership. What changes is the broader policy through which the university defines obligations, authority, procedures, assistance, and, in some cases, substantive provisions governing the relationship between inventors and the institution.

The fact that the strongest difference appears in licensing makes the organizational literature on technology transfer offices especially relevant. TTOs do much more than receive disclosures and pass inventions along to firms. Their staffing, incentives, routines, evaluation practices, and relationships with industry differ across universities \citep{ThursbyKemp2002,SiegelWaldmanLink2003,MarkmanEtAl2005,DebackereVeugelers2005}. \citet{JensenThursbyThursby2003} treat the TTO as an intermediary with objectives and information of its own, while later work places the office within a broader system of university IP management \citep{HolgerssonEtAl2019,AaboenHolgersson2019}. Licensing is precisely where many of these organizational differences become important: inventions must be evaluated, potential licensees identified, terms negotiated, and relationships with firms managed.

The lack of a similar increase in invention disclosures helps narrow the interpretation. Disclosure is closer to the point at which faculty decide whether to enter the formal technology transfer system. Previous work shows that this decision is shaped by local norms, faculty experience with the TTO, and what researchers observe among their colleagues \citep{OwenSmithPowell2001,BercovitzFeldman2008}. TTOs also need legitimacy among faculty \citep{OKaneEtAl2015}, and researchers sometimes choose informal routes when their experience with the office, outside partners, or their own commercialization expertise makes those routes more attractive \citep{HalilemDiop2025}. Perceptions of fairness appear to matter as well: \citet{WaldmanEtAl2022} find that academic entrepreneurs who view the university's commercialization process as fair are more willing to use formal channels.

If the main change associated with a policy revision were simply that more faculty decided to bring inventions into the formal system, disclosures would be a natural place to see it. We do not. Instead, the difference appears farther downstream. That is consistent with the idea that something changes in what happens after an invention has already been disclosed. It may concern evaluation, patenting decisions, marketing, negotiations with firms, or the continued relationship between the inventor and the TTO.

That last relationship is important because inventor involvement does not end with disclosure. University technologies are often licensed while still at an early stage, and continued participation by the inventor can be important for moving them toward commercial use \citep{JensenThursby2001}. Seen in that light, the positive movement in the filing margin and new patent applications, although imprecisely estimated, is also suggestive. These are all post disclosure outcomes. At the same time, our institution level data do not allow us to follow a particular invention from disclosure through patenting and licensing, so we cannot identify which part of this process is responsible for the licensing pattern.

The communication measure gives us another way to think about these results. \citet{SharmaPCSI2026} is careful to distinguish communicative stance from the substantive generosity of an IP policy. A higher PCSI does not mean that a university gives inventors better ownership rights or a larger share of licensing revenue, nor does it tell us whether faculty actually experience the TTO as supportive. It measures how the formal policy presents authority, obligations, assistance, and the relationship between the institution and its inventors.

That distinction opens several possible channels. Some faculty may encounter the policy directly when trying to understand what they must disclose or what help the university will provide. More often, the policy may matter indirectly because TTO staff work within the same formal governance framework and reproduce it through forms, guidance, training, and day to day interactions. A rewritten policy may also reflect a broader change in how university leadership wants technology transfer to operate. These possibilities are difficult to separate with our data, and they need not be mutually exclusive.

The first possibility is closely related to the literature on administrative burden. Formal procedures can impose learning and compliance costs on the people who have to navigate them \citep{MoynihanHerdHarvey2015,HerdMoynihan2018}. In principle, a policy that presents obligations and available assistance more clearly or more supportively could reduce some of those costs. But our results do not fit especially well with the simplest version of that argument, in which a more supportive policy mainly encourages additional faculty to disclose inventions. The absence of a disclosure response points us more toward changes that occur after disclosure, although it does not rule out an inventor facing mechanism altogether.

There is also a more cautious interpretation. The policy itself may not be what changes behavior. Universities rarely rewrite important governance documents in isolation. A policy revision may occur alongside new TTO leadership, different internal procedures, new staff, changes in university and industry strategy, or a broader shift in institutional priorities. \citet{MeyerRowan1977} provide a useful reminder that formal rules and organizational practice need not coincide perfectly. In our setting, a change in PCSI may therefore be one visible part of a larger institutional change rather than the independent cause of what happens to licensing.

The descriptive content results reinforce this point. Upward and downward revisions do not differ only in communicative stance. Downward revisions are more likely to add conflict of interest language, and the smaller provision coded sample shows changes in other parts of the policy as well. We do not interpret the keyword results as evidence that conflict of interest provisions themselves reduce licensing. They are useful because they show that policy rewrites often bundle together changes in language, procedure, and substantive rules. Separating those components would require independent provision level coding rather than trying to infer a mechanism from PCSI alone.

This is also why the identification problem cannot be solved only by showing similar pre revision trends. Upward and downward revisions look similar in their earlier licensing paths, and the licensing estimate remains positive under stricter sample definitions, leave one out exercises, and higher PCSI thresholds. Those checks make the pattern harder to attribute to one unusual event or to the particular 0.03 cutoff. But they do not make the direction of revision random. The most important observed difference is the predecessor policy itself: upward revisions begin from substantially lower PCSI values than downward revisions.

That difference may partly reflect simple room to move, but it may also tell us something about why universities revise. A university replacing a particularly restrictive policy may be responding to faculty dissatisfaction, a change in leadership, a TTO reorganization, or a new commercialization strategy. Any of these could affect licensing independently of the written stance of the new policy. The stacked design, contemporaneous controls, fixed effects, and research expenditure adjustment improve the comparison, but they cannot remove changes that occur together with the university's decision to revise in one direction rather than the other.

The permutation inference should be read in the same spirit. It is useful because the effective variation comes from a relatively small number of revision events, particularly only eight upward revisions. But the direction labels were not assigned by design. The permutation exercise therefore tells us how unusual the observed contrast is under an exchangeability assumption; it does not turn the analysis into a randomized experiment.

Our findings therefore sit somewhere between two parts of the existing literature. They are more modest than studies of large ownership reforms, because the universities in our sample generally remain within the same basic ownership regime. At the same time, they move beyond cross sectional comparisons of individual policy terms by following actual within university revisions and examining several stages of commercialization around those events. What emerges is a fairly specific pattern: ordinary policy revisions are associated with different licensing paths, and that difference appears more clearly after disclosure than at the point where faculty first enter the formal system.

We do not take this to mean that communication is more important than ownership, financial incentives, TTO resources, or organizational capability. A more plausible reading is that formal policy is one part of the governance relationship connecting these pieces. University IP policy tells inventors what is required of them, gives TTOs a formal framework within which to operate, and may reflect broader institutional priorities. The results suggest that this governance layer deserves more attention in explaining differences in university commercialization.

The practical implication is similarly limited. Our estimates do not imply that a university can increase licensing simply by rewriting an IP policy in more supportive language. What they do suggest is that policy revision should not automatically be treated as administrative housekeeping. In our sample, upward and downward revision episodes are followed by different licensing paths, and the difference is concentrated downstream of disclosure. For universities, the relevant object may therefore be the broader package that comes with a rewrite, including procedures, responsibilities, TTO practices, and inventor interactions, rather than the wording of the document by itself.

A useful next step would be to find settings in which universities have less choice over when or how policies change. System wide mandates, staggered implementation of a common rule across campuses, or legal changes that bind some institutions more strongly than others would provide cleaner variation. A second step would be to code the changes that accompany a rewrite in much greater detail: disclosure requirements, assignment language, appeals, conflict of interest provisions, staffing, internal review procedures, and licensing practices. Linking those changes to invention level data would make it possible to see whether the difference arises during evaluation, patenting, marketing, negotiation, or continued inventor involvement.

Several limitations remain. The preferred sample contains only eight upward revisions, and some event dates rely on tier B documentary evidence. The historical archive may also miss intermediate policy versions. AUTM reporting is voluntary, and its institution level outcomes do not allow us to follow individual inventions from disclosure to license. Finally, revisions near the end of the AUTM panel have relatively little post event support. These limitations are why we put more weight on the location and stability of the observed pattern than on giving the licensing estimate a structural causal interpretation.

\section{Conclusion}\label{sec:conclusion}

This paper shows that ordinary university IP policy revisions are followed by different licensing paths, with the clearest divergence appearing after invention disclosure rather than at the point where faculty first enter the formal technology transfer system. The contribution is not to displace existing explanations based on ownership, incentives, or TTO capacity, but to make more visible an underlying governance mechanism that is difficult to observe directly: how universities structure and communicate their relationship with inventors. Because these revisions often combine changes in language, procedures, and organizational practice, future research should look for settings with cleaner external variation and richer invention level data that can identify which parts of the commercialization process drive the observed differences.
